\clearpage
\begin{explanationbox}
\begin{minipage}{\linewidth}
\paragraph{先想象一条路：导数告诉你正在上坡还是下坡}
\begin{center}
\begin{tikzpicture}
\begin{axis}[width=11.6cm,height=6.3cm,axis lines=middle,
 xmin=-1.5,xmax=1.65,ymin=-2.5,ymax=2.7,
 xlabel={$x$},ylabel={$y$},ticks=none,clip=false,
 title={同一条曲线上：正切线斜率与高度增加},title style={font=\small}]
\addplot[blue,very thick,domain=-1.35:1.35,samples=100]{x+x^3/3};
\addplot[orange!85!black,thick,dashed,domain=-1.3:-0.3]{1.64*(x+0.8)-0.9706667};
\addplot[orange!85!black,thick,dashed,domain=0.3:1.3]{1.64*(x-0.8)+0.9706667};
\addplot[only marks,blue] coordinates {(-0.8,-0.9706667) (0.8,0.9706667)};
\node[font=\small,anchor=north west] at (axis cs:-1.2,-1.65) {$f'(-0.8)>0$};
\node[font=\small,anchor=south west] at (axis cs:0.1,2) {$f'(0.8)>0$};
\draw[densely dotted] (axis cs:-0.8,0)--(axis cs:-0.8,-0.9706667);
\draw[densely dotted] (axis cs:0.8,0)--(axis cs:0.8,0.9706667);
\end{axis}
\end{tikzpicture}
\end{center}

\end{minipage}

图中的函数是
\[
f(x)=x+\frac{x^3}{3},\qquad f'(x)=1+x^2>0.
\]
因为 $x^2\ge0$，所以 $1+x^2$ 一直大于 $0$。图中蓝线持续上升，橙色虚线是两处切线；本例整段导数均为正，由此可判断整段严格递增。

\textbf{把图形翻译成数学语言：}
\begin{itemize}
  \item $f'(x)>0$：此处曲线向右走时是向上的。
  \item $f'(x)<0$：此处曲线向右走时是向下的。
  \item $f'(x)=0$：此处切线是水平的。\textbf{仅凭这一点，不能判断附近是升、降还是掉头。}
\end{itemize}

\textbf{为什么特别强调``同一段''？}
导数告诉我们局部的变化方向；要判断从一个点走到另一个点，必须确保这两个点之间的路没有断开。分式函数常在分母为 $0$ 的位置断开，因此要分段讨论。

\textbf{一个容易记住的画面：}
``在某一点走平了''，不等于``前后有一整段平路''。例如 $y=x^3$ 在原点切线水平，左边到右边的函数值却始终在增加。反过来，若整段都是水平线，两个不同位置的函数值就会相等，不能叫严格递增。

\textbf{做题顺序：}先找能取哪些 $x$，再求 $f'(x)$，接着看正负号，最后分别写出每一段的增减情况。
\end{explanationbox}
